\subsection{N-013: descent in an explicitly bar-fixed globalization variant}
\label{sec:nonstandard-bar-fixed-descent}

\paragraph{The variant and the claim.}
This subsection changes the domain of GCT VIII, subsection 3.3,
Conjecture 3.2, by explicitly imposing $y=\bar y$.  It does not
reinterpret that condition as part of the printed conjecture, which
is refuted under its displayed domain by N-012.  We work over
$K=\mathbb Q(v)$, put $q=v^2$, $\delta=v+v^{-1}$ and $D=\delta^2$,
and allow the scaling exponents $a\in\tfrac12\mathbb Z$, precisely
the rational powers $q^a$ that belong to $K$.

Let
\[
 R=\mathbb Q[v,v^{-1}],\qquad S=\mathbb Q[D],\qquad
 F=\mathbb Q(D),\qquad
 B_S=S\langle P_1,P_2,P_3\rangle,\qquad
 B_R=R\langle P_1,P_2,P_3\rangle.
 \tag{N-44}
\]
These are the actual subalgebras generated by the indicated words
inside the Hecke tensor algebra.  No saturation or localization
of $B_S$ or $B_R$ is taken.  Let $\{c'_T\}$ be the rational GT
matrix-unit basis supplied by N-010, and fix a genuine finite
partial order on its labels.  In particular, every $c'_T$ belongs
to $B_F=F\langle P_1,P_2,P_3\rangle$.  This is a specified rational
normalization and an additional choice of label order; it does
not claim that VIII uniquely prescribes either choice.

Let $A_0\subset K$ be the ring regular at $v=0$.  At an inductive
stage $T$, suppose that all lower scaling exponents $a_U$ have
been fixed.  Define
\[
 L_a=A_0q^ac'_T\oplus
       \bigoplus_{U<T}A_0q^{a_U}c'_U,
 \qquad
 \mathcal F_a=
 \{y\in B_R\cap L_a\cap\bar L_a:
       \bar y=y,\ y-q^ac'_T\in vL_a\}.
 \tag{N-45}
\]
The degree order is exactly VIII's external order on ambient KL
coefficient pole orders at $q=\infty$; the ambient coordinate
order can be any of the finite orders described there.

The following conclusions hold.
\begin{enumerate}
\item $B_S$ is a finite free $S$-module of rank $114$, and
      $B_R=R\otimes_S B_S$ inside the actual ambient algebra.
\item A global minimum over the union of the fibres
      $\mathcal F_a$ cannot have $a\in\tfrac12+\mathbb Z$.
\item If a global minimum exists, then at the same integral
      exponent there is a minimum in $B_S$.  In particular, if
      the minimizing pair $(a,e)$ is unique, then $a\in\mathbb Z$
      and $e\in B_S$.
\end{enumerate}
These statements do not prove existence or uniqueness of a
minimum.  Their consequence is conditional descent for this
specified variant, not a construction of VIII's printed $E(4)$.

\paragraph{The actual polynomial order is free.}
The ring $R$ is finite free of rank four over $S$.  Indeed
\[
 v^4+(2-D)v^2+1=0,\qquad
 v^{-1}=(D-2)v-v^3,
 \tag{N-46}
\]
so $1,v,v^2,v^3$ span.  They are independent over $F$: the four
automorphisms $v\mapsto v,v^{-1},-v,-v^{-1}$ of $K/F$ are distinct,
whereas (N-46) bounds $[K:F]$ by four.  Thus $[K:F]=4$, proving
the asserted independence over both $F$ and $S$.

The ambient integral Hecke tensor algebra, after replacing
integer coefficients by rational coefficients, is finite free
over $R$ and hence finite over $S$.  Since $S$ is noetherian,
its submodule $B_S$ is finitely generated.  It is torsion-free,
being a submodule of an $S$-torsion-free module.  The ring $S$
is a principal ideal domain, so $B_S$ is free.  Its rank is
$\dim_F B_F=114$ by the inherited dimension and $F$-form inputs
used in N-010.

Choose an $S$-basis $b_1,\ldots,b_{114}$ of $B_S$.  Every $P$-word
is an $S$-linear combination of these basis elements, so
$R\otimes_S B_S\to B_R$ is surjective.  The $b_j$ are also
$K$-independent, because $B_F$ is an $F$-form of the actual
generic algebra.  Therefore the map is injective, establishing
the first conclusion with no change to the word order.
This freeness over $\mathbb Q[D]$ is consistent with N-009:
that result concerns the different coefficient ring
$\mathbb Z[v,v^{-1}]$ at a mixed-characteristic maximal ideal.

\paragraph{Two commuting semilinear involutions.}
Let bar be the usual Hecke involution, with $\bar v=v^{-1}$
and $\overline{T_i}=T_i^{-1}$.  Every $P_i$ and every element
of $S$ is fixed by bar.  Define a second semilinear involution
$\beta$ by
\[
             \beta(v)=-v,\qquad \beta(T_i)=-T_i.
 \tag{N-47}
\]
The quadratic relation is preserved because both $T_i$ and
$v-v^{-1}$ change sign; each braid relation is homogeneous
and has the same sign on its two sides.  Hence (N-47) defines
an involution of the Hecke algebra and its tensor product.
It commutes with bar.  Since $C'_i=T_i+v^{-1}$ and
$C_i=T_i-v$ both change sign, $P_i=C'_i\otimes C'_i+C_i\otimes C_i$
is fixed by $\beta$.  Also $\beta(D)=D$.

Consequently every $b_j$ and every $c'_T\in B_F$ is fixed
by both involutions.  In particular, $\beta$ does not permute
the GT lines or their labels.  It preserves every span cut out
by the chosen label order.  This follows from the actual
$F$-word expressions, not from an assumed symmetry of the
abstract branching diagram.

The invariant Laurent ring is $R^{\mathrm{bar}}=\mathbb Q[\delta]$:
each symmetric Laurent polynomial is a polynomial in
$v+v^{-1}$, by the recurrence
$v^{n+1}+v^{-n-1}=\delta(v^n+v^{-n})-(v^{n-1}+v^{-n+1})$.
Using the independent $b_j$ and $\beta(\delta)=-\delta$ gives
\[
 B_R^{\mathrm{bar}}=B_S\oplus\delta B_S,\qquad
 (B_R^{\mathrm{bar}})^{\beta=+1}=B_S,\qquad
 (B_R^{\mathrm{bar}})^{\beta=-1}=\delta B_S.
 \tag{N-48}
\]
These are identities in the actual order, so division by
$\delta$ on the last summand remains inside $B_R$.

\paragraph{The effect on degrees and on the fibres.}
In either usual KL convention,
$\beta(C_w)=(-1)^{\ell(w)}C_w$.  For example, this follows
from the standard expression for $C_w$ in which the exponent
of every coefficient of $T_x$ is congruent to
$\ell(w)-\ell(x)$ modulo two; replacing $v$ and every $T_i$
by their negatives gives the displayed sign.  Thus every
ambient KL tensor coordinate of $\beta y$ is the corresponding
coordinate of $y$, evaluated at $-v$ and multiplied by a sign.
It follows that
\[
 d(\beta y,z)=d(y,z),\qquad
 d(y_1+y_2,z)\le\max\{d(y_1,z),d(y_2,z)\}.
 \tag{N-49}
\]
The convention $d(0,z)=-\infty$ is used here and below.

Fix $a\in\tfrac12\mathbb Z$ and set $s=(-1)^{2a}$.
Each generator $q^{a_U}c'_U$ of the lower lattice and
$q^ac'_T$ is an eigenvector of $\beta$ with eigenvalue $+1$
or $-1$.  The coefficient ring $A_0$ is preserved by $\beta$,
and its residue field is $\mathbb Q$, on which $\beta$ is the
identity.  Therefore, if $y\in\mathcal F_a$, then
$s\beta(y)\in\mathcal F_a$.  The vector
\[
                   y_s=\tfrac12(y+s\beta(y))
 \tag{N-50}
\]
also belongs to $\mathcal F_a$, is nonzero by its specified
residue, and satisfies $\beta(y_s)=s y_s$.
Its degrees are coordinatewise at most those of $y$ by
(N-49).  Averaging preserves membership in the actual order
because its coefficient ring is rational; no division by two
in the integral order of N-009 is asserted.

If $s=-1$, (N-48) permits $z=y_s/\delta\in B_S$.
In the $T$ coordinate, division by $\delta$ changes the
admissible exponent from $a$ to $a+1/2$, since
\[
 \frac{q^a/\delta}{q^{a+1/2}}=\frac1{1+q}\in1+vA_0.
 \tag{N-51}
\]
In every lower coordinate, the factor
$1/\delta=v/(1+v^2)$ increases its valuation at zero by one
in the $v$ variable.  A coefficient already in $vA_0$
therefore remains in $vA_0$.  Hence $z\in L_{a+1/2}$ has the
required residue.  It is bar fixed, so membership in
$\bar L_{a+1/2}$ follows by applying bar to its membership
in $L_{a+1/2}$.  Thus
\[
                 y_s/\delta\in\mathcal F_{a+1/2}.
 \tag{N-52}
\]
At every coordinate where $y_s$ is nonzero, division by
$\delta$ decreases its degree at $q=\infty$ by $1/2$.
At every coordinate where $y$ is nonzero but $y_s$ is zero,
the new degree is $-\infty$.  Therefore (N-49)--(N-52)
strictly decrease the degree at every nonzero coordinate of
$y$, and never introduce a new nonzero coordinate.

This gives strict decrease in VIII's partial degree order,
not merely in a total refinement.  At a minimal coordinate
of the nonzero support of $y$, the reverse comparison fails:
the new degree is strictly smaller and all preceding
coordinates have value $-\infty$ for both vectors.
Thus no candidate with $s=-1$ can be a global minimum.

If $s=+1$ and $y$ is a global minimum, (N-50) gives a vector
$y_+\in B_S\cap\mathcal F_a$ with coordinatewise no larger
degree.  It is also a global minimum: for any competing
candidate $w$, each comparison witnessing $d(y)\le d(w)$
still holds with $y_+$ in place of $y$.  Indeed a coordinate
inequality remains true after lowering that coordinate, and
a strictly smaller preceding coordinate remains strictly
smaller after such lowering.  If the global minimum is
unique, $y=y_+\in B_S$.  This proves all three assertions.

\paragraph{What the descent does and does not settle.}
For the stated rational normalization, a unique bar-fixed
external globalization, if it exists, does not require a
coefficient extension of $\mathbb Q(D)$ at rank four.
This conclusion is stronger than merely splitting the generic
algebra: the resulting vectors would lie in its actual
$\mathbb Q[D]$ word order.  Their multiplication coefficients
in that rational basis would then belong to $F$, although
finite Laurent-polynomial or signed-positivity properties
still require separate proofs, as N-007 demonstrates.

There is also a concrete finite-linear-algebra consequence.
Use the flip-symmetric ambient KL coordinates and write a
$\beta$-fixed, bar-fixed coordinate as
$\delta^{\epsilon_z}p_z(D)$, where
$\epsilon_z\equiv\ell(w_1)+\ell(w_2)\pmod2$ and
$p_z\in\mathbb Q[D]$.  A fixed finite upper bound $h_z$
on every nonzero ambient degree permits only
$\deg p_z\le h_z-\epsilon_z/2$.  Within that finite-dimensional
space, actual-order membership, the lower-label support
conditions, and any fixed-scaling residue conditions are
linear conditions over $\mathbb Q$ once a polynomial module
basis is supplied.  This yields an exact feasibility problem
for every bounded box.

It does not, by itself, replace the unbounded optimization by
one bounded box.  A strict improvement at an earlier ambient
coordinate can permit arbitrarily large later coordinates in
VIII's partial degree order.  Nor does it select an order for
modules whose maximal $H$-weights are incomparable.  These
existence, uniqueness, order and global-cutoff issues remain
explicit gates for the repaired construction.
